\documentclass[11pt,a4paper]{article}
\usepackage[T1]{fontenc}
\usepackage[utf8]{inputenc}
\usepackage{lmodern}
\usepackage[margin=26mm]{geometry}
\usepackage{booktabs,tabularx,array}
\usepackage[colorlinks=true,urlcolor=blue,linkcolor=blue,citecolor=blue]{hyperref}
\pdfinfoomitdate=1
\pdftrailerid{}
\pdfsuppressptexinfo=15
\hypersetup{pdftitle={TOS: Post-Quantum Authentication and Shielded Transfers},pdfauthor={TOS contributors}}
\setlength{\emergencystretch}{3em}
\setlength{\parindent}{0pt}
\setlength{\parskip}{6pt}
\newcommand{\code}[1]{\texttt{#1}}
\newcommand{\repo}[2]{\href{https://github.com/tosnetwork/tos/blob/main/#1}{#2}}
\title{TOS Blockchain\\\large Post-Quantum Authentication and Shielded Transfers}
\author{TOS contributors}
\date{October 2026}
\begin{document}
\maketitle
\begin{abstract}
TOS is a Layer-1 blockchain with masterchain and shardchain infrastructure,
asynchronous smart contracts, and a native virtual machine. Its canonical
genesis configuration uses ML-DSA-44 validator signatures. Its shielded
native-asset design combines zero-knowledge transaction proofs, one-time
post-quantum note authorization, and post-quantum encrypted note delivery.
These mechanisms protect different surfaces: signature security is not the
same property as transaction privacy or proof soundness. In particular, the
current shielded circuit uses pairing-based Groth16 proofs and is not a
post-quantum proof system. This overview describes implementation boundaries,
public information, and deployment prerequisites rather than asserting
unconditional anonymity or production readiness.
\end{abstract}
\textbf{Scope.} Source baseline: \code{54be8f1f4}. This document describes the
repository's October 2026 implementation and canonical configuration; it is
not a statement that every running network has activated them. Exact wire
formats and deployment identities remain defined by source and the frozen
shielded-pool profile.
\setcounter{tocdepth}{1}
\tableofcontents
\newpage

\section{Architecture and security objectives}
The masterchain carries network configuration and validator-set information.
Shardchains execute accounts and transfer messages. The TOS Virtual Machine
(TVM) executes contracts, while validators authenticate consensus messages.
Contract calls between accounts are asynchronous messages: receiving funds,
forwarding them, and processing a later result are separate transactions.
A signature cannot by itself guarantee completion or refund of that sequence.

TOS applies cryptographic mechanisms at distinct layers:
\begin{center}
\small
\begin{tabularx}{\linewidth}{@{}p{0.23\linewidth}X@{}}
\toprule
Surface & Mechanism and boundary \\
\midrule
Validator signing & ML-DSA-44; protects authorized consensus messages. \\
Contract authorization & Native ML-DSA-44 verification, with contract-specific signed fields and replay rules. \\
Shielded note spending & One-time ML-DSA-44 note authorization plus a circuit proof of valid state transition. \\
Note delivery & ML-KEM-768 key encapsulation and XChaCha20-Poly1305 authenticated encryption. \\
Hidden-value validity & Groth16 over BLS12-381, with Poseidon2 commitments and tree hashing; not post-quantum proof soundness. \\
\bottomrule
\end{tabularx}
\end{center}
The objective is to provide post-quantum authentication on designated paths
and confidentiality within the shielded protocol. It is not to imply that
all legacy wallets, administrative credentials, transport authentication, or
external applications become post-quantum merely by connecting to TOS.

\section{Post-quantum consensus from genesis}
\subsection{Canonical version parameters}
The canonical zerostate generator selects global protocol version 16 in
ConfigParam 8, and creates four equally weighted post-quantum genesis
validator descriptors. These use \code{validator\_pq\_addr\#b3} and algorithm
identifier 1 for ML-DSA-44. ConfigParam 47 binds controller admission to the
approved controller code. The Simplex transport configuration selects its
v2 QUIC profile; that transport version is distinct from global protocol
version 16.

The public bootstrap manifest is \code{validator-pq.pub}. It binds each
validator's controller identity, ADNL identity, and consensus public key.
The canonical generator rejects the classical bootstrap-file format and
checks duplicate identities. Private keys remain operator material and must
not be placed in the public manifest or committed to the repository.
Launching a new chain still requires review of all immutable genesis inputs,
including network identity, time parameters, and compiled contract hashes.
Changing the generator does not change an existing chain's zerostate.

\subsection{Keys and authorization}
ML-DSA is standardized in FIPS 204~\cite{fips204}. The implemented ML-DSA-44
suite uses a 1,312-byte public key and a 2,420-byte signature. Protocol-specific
message encodings and domain separation remain essential: algorithm security
does not make a signature valid for another role, network, or request.

A stable controller identity separates validator identity from individual
keys. Root authority and consensus signing have different responsibilities;
ADNL identity is also distinct. Rotation and admission follow the controller
and network rules, rather than an assumption that every address or key is
interchangeable. Operators must provision, protect, and recover the correct
key for each role.

TVM exposes \code{PQCHECKSIG\_MLDSA44} for contract verification. Its minimum
global version is 16; the canonical genesis profile is version 16. An
instruction's introduction version is not the complete version requirement
of a contract using additional instructions.

\subsection{Resource implications}
Post-quantum signatures increase serialized message sizes relative to compact
classical signatures. Admission limits, verification budgets, message-size
limits, and network capacity must account for those sizes. This overview
makes no throughput or fee claim: measured cost depends on the implementation,
message path, hardware, and workload. Native contract verification also does
not imply that each contract call repeats consensus verification.

\section{Shielded native-asset transfers}
\subsection{State and transaction flow}
The version-1 shielded pool targets native TOS in workchain 0. It records
commitments and nullifiers rather than exposing a transparent account balance
for each private note. Commitments bind note contents; nullifiers prevent a
consumed note from being spent again. Proof verification checks the specified
transition, including the circuit's value and membership constraints.

\begin{center}
\fbox{\begin{minipage}{0.89\linewidth}
Transparent TOS deposit $\longrightarrow$ shielded notes\\[4pt]
Shielded notes $\longrightarrow$ proof-authorized private transfer
$\longrightarrow$ new notes\\[4pt]
Shielded notes $\longrightarrow$ public withdrawal message
\end{minipage}}
\end{center}
Each private transaction has a fixed shape of two input slots and three
output slots, with real and dummy data handled according to the profile.
The circuit uses Groth16 on BLS12-381, and Poseidon2 over its scalar field
for commitments and tree operations. Tree parameters, public inputs, proof
encoding, and deployment hashes must match the frozen profile exactly.

Note authorization uses one-time ML-DSA-44 keys. Authorization and proof
verification serve different purposes: the former authenticates spending
intent, while the latter establishes the hidden-state relationships required
by the circuit. Neither can substitute for the other.

\subsection{Encrypted delivery}
Recipients discover and decrypt note payloads using the wallet's viewing
material. Delivery combines ML-KEM-768, standardized in FIPS 203~\cite{fips203},
with XChaCha20-Poly1305. A domain-separated derivation binds key material to
the execution domain and output slot, with authenticated context preventing
payloads from being silently reassigned to another context.

The fixed payload contains a version byte, a 1,088-byte encapsulation
ciphertext, and a 144-byte authenticated encrypted body: 1,233 bytes in total.
Real and dummy outputs use the same payload size. This removes a simple
length distinction, not all statistical or network side channels.

Payment descriptors are single-use. Viewing keys can remain shared within a
pool domain; senders who exchange off-chain descriptor information can
recognize a reused viewing key. The current profile therefore does not claim
unlinkability against such sender collusion.

\subsection{What remains public}
Shielded transfers do not hide every observable fact. Transparent deposit
values, public withdrawal values and recipients, fees, commitments,
nullifiers, transaction ordering, and timing remain observable as specified
by the protocol. A node operator or network observer may also see traffic
metadata. Wallet compromise can reveal note or viewing material.

Privacy claims must consequently identify the observer and the protected
information. Hidden note values and ownership relationships are not the same
as hiding network origin, deposit history, or every relationship inferable
from user behavior. Wallet practices and the size and activity of the privacy
set affect practical privacy.

\section{Cryptographic and financial boundaries}
\subsection{Post-quantum scope}
ML-DSA signatures and ML-KEM delivery address quantum threats to their
respective authentication and encapsulation tasks. Groth16 is a pairing-based
proof construction~\cite{groth16}; its soundness assumptions are not
post-quantum. Adding a post-quantum signature to a Groth16 transaction does
not make the proof system post-quantum.

The frozen profile explicitly defers a post-quantum proof backend. Future
migration would require new proof and circuit review, compatible state and
wallet handling, and a controlled activation process. It cannot be represented
as an already-delivered property of the current pool.

\subsection{Funds, fees, and asynchronous outcomes}
A valid proof must still be processed by a correctly funded contract. The
pool must preserve backing and enforce fee and message rules. Deposit and
transfer entry points use funded internal messages; proof acceptance does
not authorize arbitrary expenditure of the account's balance.

For withdrawal, the public amount is the gross outgoing message value before
recipient computation, not a guarantee of the recipient's net balance
increase. Failure and bounce handling must account for the value actually
returned. Fees already consumed by execution or delivery cannot be assumed
to reappear as refundable principal. This is an asynchronous settlement
boundary, independent of whether the authorization is post-quantum.

The v1 profile does not provide arbitrary multi-asset transfers, variable
transaction shapes, or unrestricted payload memos. Its state structures and
history are bounded according to the profile. These bounds require explicit
capacity and recovery testing; the word ``shielded'' is not an exemption
from storage or execution constraints.

\section{Implementation and deployment status}
The repository contains the pool contract, circuit, wallet implementation,
and deployment artifacts. Availability of those files is not activation
approval. The frozen profile labels itself ready for coding and requires its
acceptance gates before mainnet or testnet activation.

\textbf{Development proving keys must never protect real funds.} The circuit
package describes deterministic development setup material whose secret
setup randomness is known. Such material allows forged proofs. A production
deployment requires the prescribed ceremony and independently verified
artifacts, not merely rebuilding the development example.

Before activation, operators and reviewers must establish at least:
\begin{enumerate}
\item Agreement among contract code, circuit, public-input encoding, verifier
keys, and deployment-profile identities.
\item Completion of the required ceremony and verification of its final
artifacts, without substituting development keys.
\item Native execution of valid and invalid paths, including replay, double
spending, malformed payloads, insufficient funding, and asynchronous failures.
\item Resource and recovery checks for the chosen network configuration.
\item Explicit network activation approval under the release and upgrade
process; no assumption that source availability authorizes deployment.
\end{enumerate}
This is a summary, not a replacement for the profile's acceptance matrix.
There is no administrative \code{SETCODE} upgrade path in the pool design;
changes require the specified network upgrade process.

\section{Development and operations}
The repository supplies the validator engine, clients, TVM, contract tools,
and supporting packages. See \repo{BUILD.md}{the build guide} for dependencies
and target commands. \repo{doc/Local-PQ-Network.md}{The local PQ network guide}
covers local deployment and election rehearsals; local test success is not a
production security certification.

Monitoring should distinguish unavailable observations from a healthy node.
Process liveness alone does not establish validator participation, progress,
or storage health. Cryptographic validation, operational observations, and
fund-recovery tests should each report their own boundaries.

This overview is built from \code{doc/pq.tex} using
\code{bash scripts/build-pq-paper.sh}. The script requires \code{pdflatex}
and writes \code{doc/pq.pdf} without leaving intermediate files in the
source directory. Historical protocol papers remain available separately.

\section*{Implementation references}
\addcontentsline{toc}{section}{Implementation references}
\begin{itemize}
\item \repo{crypto/smartcont/gen-zerostate.fif}{Canonical genesis generator}
and \repo{doc/validator-genesis-bootstrap.md}{validator bootstrap guide}.
\item \repo{crypto/pq/pq-consensus.h}{Post-quantum consensus types} and
\repo{doc/tvm-mldsa44.md}{native ML-DSA verification interface}.
\item \repo{artifacts/shielded-pool/PROFILE.md}{Frozen shielded-pool profile}:
exact encoding, privacy boundaries, and activation gates.
\item \repo{tools/shielded-pool-circuit/README.md}{Circuit package and development-key warning}.
\item \href{https://github.com/tosnetwork/tos/tree/main/tools/shielded-pool-wallet}{Shielded wallet implementation}.
\item \repo{artifacts/shielded-pool/CEREMONY.md}{Ceremony requirements} and
\repo{doc/tos-upgrade-process.md}{network upgrade process}.
\end{itemize}
\begin{thebibliography}{9}
\bibitem{fips204} NIST. \emph{Module-Lattice-Based Digital Signature Standard}.
FIPS 204, 2024. \url{https://doi.org/10.6028/NIST.FIPS.204}.
\bibitem{fips203} NIST. \emph{Module-Lattice-Based Key-Encapsulation Mechanism Standard}.
FIPS 203, 2024. \url{https://doi.org/10.6028/NIST.FIPS.203}.
\bibitem{groth16} Jens Groth. \emph{On the Size of Pairing-Based Non-interactive
Arguments}. EUROCRYPT, 2016.
\url{https://doi.org/10.1007/978-3-662-49896-5_11}.
\end{thebibliography}
\end{document}
